Křivka prvotní magnetizace $M(H)$ začíná v odmagnetovaném stavu a má čtyři oblasti:
\begin{enumerate}
\item \textbf{I: vratné posuny Blochových stěn}; domény orientované výhodně vůči poli se zvětšují, magnetizace roste mírně.
\item \textbf{II: nevratné posuny stěn}; výhodně orientované domény rostou na úkor ostatních, magnetizace roste strmě.
\item \textbf{III: stáčení vektoru spontánní magnetizace} do směru vnějšího pole.
\item \textbf{IV: nasycení a paraproces}; směr magnetizace se již nemění, její velikost při dalším zesilování pole jen nepatrně roste.
\end{enumerate}
